% Generated draft: requires mathematical and source review.
\documentclass[11pt]{article}
\usepackage[margin=25mm]{geometry}
\usepackage{fontspec}
\setmainfont{Latin Modern Roman}
\usepackage{amsmath,amssymb,mathtools,unicode-math}
\setmathfont{Latin Modern Math}
\usepackage{enumitem,needspace,xurl,hyperref}
\setlength{\parindent}{0pt}
\setlength{\parskip}{5pt}
\setlength{\emergencystretch}{3em}
\newcommand{\N}{\mathbb{N}}
\newcommand{\Z}{\mathbb{Z}}
\newcommand{\Q}{\mathbb{Q}}
\newcommand{\R}{\mathbb{R}}
\newcommand{\C}{\mathbb{C}}
\newcommand{\F}{\mathbb{F}}
\newcommand{\A}{\mathbb{A}}
\newcommand{\PP}{\mathbb{P}}
\newcommand{\E}{\mathbb{E}}
\newcommand{\eps}{\varepsilon}
\newcommand{\dd}{\,\mathrm{d}}
\begin{document}
\section*{UnsolvedMath ID: 20000348}
\textbf{Source problem code:} catalogue entry
\subsection{Definitions and mathematical statement}
Let \(A_g\) be the complex moduli space of principally polarized abelian varieties of dimension \(g\), and let \(A_g^F\) denote its first Voronoi, or perfect-cone, compactification. Write \(M\) for the weight-one modular-form or Hodge class and \(D\) for the boundary divisor; on a coarse space these may require \(\Q\)-Cartier interpretations. Shepherd-Barron's source problem has six parts.
\begin{enumerate}
\item Run the MMP on \(A_g^F\) for \(g\leq10\), particularly flipping its stated extremal ray. The source additionally says \(A_g^F\) is the canonical model of \(A_g\) for \(g\geq12\), and a minimal model for \(g=11\); these are supplied background claims.
\item For the natural inclusions \(A_g^F\to A_{g+1}^F\) parameterized by the elliptic \(j\)-line, ask whether \(H^*(A_g^F,\C)\) stabilizes and identify the limit. A precise formulation needs degreewise ranges and the induced pullback maps, rather than unstated stabilization of the entire finite-dimensional graded group at once.
\item Compute the intersections \(M^aD^b\) in top degree \(a+b=g(g+1)/2\), and the plurigenera of \(A_g^F\), with coarse/stack and canonical-power conventions fixed.
\item Determine whether its total coordinate ring is finitely generated.
\item Determine whether \(12M-D\) is semiample for \(g\geq12\), meaning an appropriate positive Cartier multiple is globally generated.
\item Find its effective cone, related in the source to possible slopes of Siegel cusp forms.
\end{enumerate}

For \(g\geq2\), the attached report proposes a bridge between the negative-boundary ring and a lattice Cox ring after finite-index Cartier regrading. A common interpretation, once \(NM\) and \(ND\) are Cartier, is the section ring with degrees \(aNM-bND\), \(a,b\geq0\), compared with the multisection ring for the lattice \(\Lambda=\Z(NM)+\Z(ND)\). The report asserts that adjoining the canonical section of the effective Cartier boundary multiple generates that lattice Cox ring from the exact Shepherd-Barron negative-boundary bigraded ring. The precise definition of that ``exact'' ring, grading lattice, and relation to a full Weil-class Cox ring are not printed in the summary; these must be checked, rather than silently replacing it by a different ring.

The report claims that finite generation of this ring makes the compactification a Mori dream space, forces \(12M-D\) to be semiample in genus \(g\) and every lower genus, and makes the effective slope rational and attained. The usual effective slope is the infimum of \(a/b\) over effective classes \(aM-bD\) with \(b>0\), in the specified modular-form convention. The Mori-dream-space conclusion needs particular care: the standard definition includes \(\Q\)-factoriality and other Picard hypotheses, and the source does not provide these for every perfect-cone compactification or specify a \(\Q\)-factorial model. This remains a reported implication with a missing geometric qualification, not an unconditional verified theorem.

The report further claims a fan-level primitive-ray proof of compatibility under a fixed-elliptic product inclusion \(i_E:A_h^F\to A_{h+1}^F\). In the indicated Hodge/boundary normalization the asserted restriction behavior is
\[
i_E^*M_{h+1}=M_h,\qquad i_E^*D_{h+1}=D_h,\qquad
i_E^*(12M_{h+1}-D_{h+1})=12M_h-D_h.
\]
It uses this to propagate semiampleness to lower genus and says it obtains explicit coarse and stack product-locus intersection recurrences. Those recurrence formulas and their stabilizer/multiplicity normalizations are absent from the record. The candidate does not run the requested low-genus MMP, compute all plurigenera or intersections, or determine the stable cohomology.
\subsection{Short English statement}
The source asks six questions about the perfect-cone compactification: its MMP, stable cohomology, intersections, coordinate ring, semiampleness, and effective cone. The report proposes a finite-generation bridge with lower-genus restrictions, while leaving ring conventions and Mori-dream-space hypotheses unresolved.
\subsection{Sources}
\begin{itemize}
\item AIM workshop ``Compact moduli spaces and birational geometry,'' Shepherd-Barron's question extracted as ``Problem 6. 2.'': \url{https://aimath.org/WWN/birational/birational.pdf}. No PDF was accessed here.
\item Local canonical registry \texttt{aim-algebraic-geometry-notes.json}, \texttt{notes[347]}, tag \texttt{problem}, source item recorded as 6; permanent UnsolvedMath ID \(20000348\), source code \texttt{catalogue entry}. The source-verification field says the OCR string \texttt{AFg} was recovered as \(A_g^F\) from \url{https://aimath.org/WWN/birational/birational.tex}, lines 555--572 at its historical access, and checked against the PDF. This supplied verification claim is recorded without accessing either source here.
\item Supplied research attempt 1 proposing the Cox-restriction bridge and primitive-ray/product-locus calculations. No full authoritative bibliographic references, precise ring definition, or recurrence formulas are supplied. The report is machine generated, not peer reviewed, and requires independent verification before citation.
\end{itemize}
\end{document}
